%!TEX program = xelatex
\documentclass[aspectratio=169]{ctexbeamer}
\usepackage[bluetheme]{ustcbeamer}
\input{ustctheme.tex}

\usepackage{listings}
\usepackage{booktabs}
\usepackage{multirow}
\usepackage{array}
\usepackage{textcomp}

\definecolor{codebg}{RGB}{245,245,245}
\definecolor{codeframe}{RGB}{180,180,180}
\definecolor{codekw}{RGB}{0,0,180}
\definecolor{codecm}{RGB}{100,100,100}

\lstset{
  basicstyle=\ttfamily\scriptsize,
  keywordstyle=\color{codekw}\bfseries,
  commentstyle=\color{codecm},
  stringstyle=\color{red},
  numbers=left,
  numberstyle=\tiny\color{gray},
  frame=single,
  backgroundcolor=\color{codebg},
  rulecolor=\color{codeframe},
  breaklines=true,
  language=C,
  showstringspaces=false,
  tabsize=2,
  xleftmargin=1.5em,
  framexleftmargin=1em
}

\title[9P-Walkers 项目答辩]{
  9P-Walkers\\
  基于 mini9P 的 MCU 集群分布式系统\\
  项目答辩
}
\author[9P-Walkers]{
  报告人：霍斌\\
  项目组成员：董宇皓、霍斌、刘亦航、刘子源、武文韬
}
\institute[USTC]{
  中国科学技术大学\\
  OSH-2026 课程项目
}
\date{2026年7月1日}

\begin{document}

\maketitleframe

\begin{frame}
  \frametitle{大纲}
  \tableofcontents[hideallsubsections]
\end{frame}

\AtBeginSection[]{
\setbeamertemplate{footline}[footlineoff]
  \begin{frame}
    \frametitle{大纲}
    \tableofcontents[currentsection,subsectionstyle=show/show/hide]
  \end{frame}
\setbeamertemplate{footline}[footlineon]
}

%==============================================================================
\section{项目概述}
%==============================================================================

\begin{frame}
  \frametitle{9P-Walkers 是什么}
  \begin{columns}[T]
    \begin{column}{0.55\textwidth}
      \begin{block}{核心目标}
        \small
        在异构 MCU 集群上构建一个\alert{统一命名、自动组网、支持远程调用和分布式计算}的实验性分布式系统。
      \end{block}
      \begin{block}{系统组成}
        \small
        \begin{itemize}
          \item \textbf{主机}：ESP32-P4（Ethernet）/ ESP32-S3（WiFi）
          \item \textbf{节点}：STM32F407 / STM32F429
          \item \textbf{互联}：UART 多跳有线网络 + 主机间局域网
          \item \textbf{接口}：统一路径 \texttt{/mcuN/...}
        \end{itemize}
      \end{block}
    \end{column}
    \begin{column}{0.42\textwidth}
      \begin{figure}
        \centering
        \fbox{\parbox{0.9\textwidth}{\centering\vspace{2cm}系统实物拓扑图\\（请替换为实际照片）\vspace{2cm}}}
        \caption{ESP32 与 STM32 组网实物}
      \end{figure}
    \end{column}
  \end{columns}
\end{frame}

\begin{frame}
  \frametitle{应用场景}
  \begin{columns}[T]
    \begin{column}{0.24\textwidth}
      \begin{block}{\centering 边缘计算}
        \centering
        \small
        把计算任务分发到多个 STM32 节点并行执行
      \end{block}
    \end{column}
    \begin{column}{0.24\textwidth}
      \begin{block}{\centering 远程诊断}
        \centering
        \small
        通过统一路径读取任意节点的健康状态和运行指标
      \end{block}
    \end{column}
    \begin{column}{0.24\textwidth}
      \begin{block}{\centering 分布式渲染}
        \centering
        \small
        把光线追踪 tile 任务分派到多个 MCU，结果汇总显示
      \end{block}
    \end{column}
    \begin{column}{0.24\textwidth}
      \begin{block}{\centering 多主机协同}
        \centering
        \small
        多台 ESP32 在局域网内自动发现并共享全局节点视图
      \end{block}
    \end{column}
  \end{columns}
  \vspace{0.4cm}
  \begin{figure}
    \centering
    \fbox{\parbox{0.85\textwidth}{\centering\vspace{1.2cm}应用场景示意图：机器人/传感器网络/边缘计算集群（请替换为实际图片）\vspace{1.2cm}}}
    \caption{典型应用场景}
  \end{figure}
\end{frame}

\begin{frame}
  \frametitle{要解决的五个核心问题}
  \begin{enumerate}
    \item \textbf{统一命名与访问}：如何让上层用户像访问本地文件一样访问远端 MCU 资源？
    \item \textbf{自动发现与多跳路由}：节点上电后如何自动加入网络、获取地址、找到路径？
    \item \textbf{控制面与数据面隔离}：文件访问、远程调用和计算任务如何保证不阻塞链路维护？
    \item \textbf{主机间协作}：多台 ESP32 如何形成一致的全局视图并协同管理各自子树？
    \item \textbf{长计算不阻塞通信}：资源受限的 MCU 上如何安全地执行耗时计算任务？
  \end{enumerate}
  \vspace{0.3cm}
  \begin{alertblock}{解决思路}
    \centering
    统一文件命名 + 自定义链路协议 + 分层任务调度 + 主控协调 + 增量计算执行
  \end{alertblock}
\end{frame}

%==============================================================================
\section{整体架构}
%==============================================================================

\begin{frame}
  \frametitle{两个网络平面}
  \begin{figure}
    \centering
    \fbox{\parbox{0.9\textwidth}{\centering\vspace{2.5cm}双平面架构图\\主机间 IP 平面 + STM32 有线 mesh 平面\\（请替换为 architecture.pdf/png）\vspace{2.5cm}}}
    \caption{系统双平面架构}
  \end{figure}
  \begin{itemize}
    \item \textbf{主机间 IP 平面}：ESP32-P4（Ethernet）与 ESP32-S3（WiFi）通过 mDNS + TCP/9909 + CBOR 互联。
    \item \textbf{STM32 有线 mesh 平面}：节点通过 UART 多跳连接，主机负责地址分配和路由下发。
    \item 两个平面只在 ESP32 主机汇合，跨主机访问由 owner 主机代理。
  \end{itemize}
\end{frame}

\begin{frame}
  \frametitle{主机侧架构}
  \begin{columns}[T]
    \begin{column}{0.48\textwidth}
      \begin{block}{核心模块}
        \small
        \begin{itemize}
          \item \texttt{coordinator\_runtime}：UART RX 唯一消费者
          \item \texttt{host\_coordinator}：地址、租约、拓扑、路由
          \item \texttt{session\_manager}：并发 pending 管理
          \item \texttt{cluster\_vfs}：\texttt{/mcuN/...} 命名空间
          \item \texttt{slave\_rpc / job\_manager}：远程调用与任务
          \item \texttt{host\_rpc}：主机间发现与协作
          \item \texttt{web / shell / lua}：用户入口
        \end{itemize}
      \end{block}
    \end{column}
    \begin{column}{0.48\textwidth}
      \begin{figure}
        \centering
        \fbox{\parbox{0.85\textwidth}{\centering\vspace{2cm}主机模块分层图\\（请替换为 master\_arch.pdf）\vspace{2cm}}}
        \caption{ESP32 主机内部模块}
      \end{figure}
    \end{column}
  \end{columns}
\end{frame}

\begin{frame}
  \frametitle{节点侧架构}
  \begin{columns}[T]
    \begin{column}{0.48\textwidth}
      \begin{block}{分层结构}
        \small
        \begin{itemize}
          \item UART DMA / IDLE 中断
          \item 端口管理：HELLO 发现、上游选择
          \item 节点控制：地址、租约、路由、转发
          \item 服务运行时：mini9P server、RPC、Job
          \item 本地 VFS：\texttt{/sys/*}、\texttt{/compute/*}
          \item 计算 worker：增量执行 kernel
        \end{itemize}
      \end{block}
    \end{column}
    \begin{column}{0.48\textwidth}
      \begin{figure}
        \centering
        \fbox{\parbox{0.85\textwidth}{\centering\vspace{2cm}STM32 节点模块分层图\\（请替换为 slave\_arch.pdf）\vspace{2cm}}}
        \caption{STM32 节点内部模块}
      \end{figure}
    \end{column}
  \end{columns}
\end{frame}

\begin{frame}
  \frametitle{数据面类型}
  \begin{table}
    \centering
    \small
    \begin{tabular}{|c|c|c|c|}
      \hline
      \textbf{类型} & \textbf{用途} & \textbf{主机端} & \textbf{STM32 端} \\
      \hline
      mini9P & 文件/诊断 & cluster\_vfs + session\_manager & mini9P server + local\_vfs \\
      \hline
      RPC & 短调用/流 & slave\_rpc & rpc\_service \\
      \hline
      Job & 异步计算 & job\_manager & job\_service + compute\_worker \\
      \hline
      Bulk & 大数据预留 & 未实现 & 未实现 \\
      \hline
    \end{tabular}
    \caption{三类数据面对应模块}
  \end{table}
  \begin{figure}
    \centering
    \fbox{\parbox{0.85\textwidth}{\centering\vspace{1cm}三种数据面交互示意图（请替换为 dataplane.pdf）\vspace{1cm}}}
    \caption{数据面分类}
  \end{figure}
\end{frame}

%==============================================================================
\section{核心运行流程}
%==============================================================================

\begin{frame}
  \frametitle{系统启动流程}
  \begin{columns}[T]
    \begin{column}{0.48\textwidth}
      \begin{block}{ESP32-P4 主机}
        \small
        \begin{enumerate}
          \item \texttt{app\_main()} 打印芯片信息
          \item 启动 \texttt{coordinator\_runtime}
          \item 创建主循环任务和健康探测任务
          \item 启动 LAN、host RPC、WebShell、Lua、推理
        \end{enumerate}
      \end{block}
    \end{column}
    \begin{column}{0.48\textwidth}
      \begin{block}{STM32F407 节点}
        \small
        \begin{enumerate}
          \item \texttt{main()} 初始化 HAL 与外设
          \item FreeRTOS 初始化
          \item 创建队列和 8 个静态任务
          \item 启动调度器
        \end{enumerate}
      \end{block}
    \end{column}
  \end{columns}
  \vspace{0.3cm}
  \begin{figure}
    \centering
    \fbox{\parbox{0.85\textwidth}{\centering\vspace{1cm}启动时序图（请替换为 boot\_sequence.pdf）\vspace{1cm}}}
    \caption{主机与节点启动流程}
  \end{figure}
\end{frame}

\begin{frame}
  \frametitle{节点上线流程}
  \begin{figure}
    \centering
    \fbox{\parbox{0.9\textwidth}{\centering\vspace{2cm}节点上线状态转换图\\HELLO $\rightarrow$ REGISTER $\rightarrow$ ASSIGN $\rightarrow$ RENEW $\rightarrow$ LINK STATE $\rightarrow$ ROUTE UPDATE\\（请替换为 node\_join.pdf）\vspace{2cm}}}
    \caption{STM32 节点上线状态机}
  \end{figure}
  \begin{enumerate}
    \item HELLO/ACK 发现邻居
    \item 选择上游端口（直连主机或可达 relay）
    \item 发送注册请求
    \item 主机分配地址和租约
    \item 周期续租
    \item 上报链路状态，接收路由更新
  \end{enumerate}
\end{frame}

\begin{frame}
  \frametitle{数据面请求流程示例}
  \begin{columns}[T]
    \begin{column}{0.48\textwidth}
      \begin{block}{\small \texttt{cat /mcu1/sys/health}}
        \small
        \begin{enumerate}
          \item 解析路径 $\rightarrow$ 节点地址
          \item 建立/复用 mini9P session
          \item attach $\rightarrow$ walk $\rightarrow$ open $\rightarrow$ read
          \item 节点 local\_vfs 生成健康文本
          \item 响应沿原路返回
        \end{enumerate}
      \end{block}
    \end{column}
    \begin{column}{0.48\textwidth}
      \begin{figure}
        \centering
        \fbox{\parbox{0.85\textwidth}{\centering\vspace{2cm}请求响应序列图\\（请替换为 request\_sequence.pdf）\vspace{2cm}}}
        \caption{读取远端节点健康状态}
      \end{figure}
    \end{column}
  \end{columns}
\end{frame}

\begin{frame}
  \frametitle{Job 提交流程}
  \begin{columns}[T]
    \begin{column}{0.48\textwidth}
      \begin{block}{流程}
        \small
        \begin{enumerate}
          \item 用户提交：\texttt{job submit mcu1 hash hello}
          \item job\_manager 分配 host\_job\_id
          \item 发送任务到 STM32
          \item job\_service 提交给 compute\_worker
          \item compute\_task 增量执行
          \item 用户查询状态和结果
        \end{enumerate}
      \end{block}
    \end{column}
    \begin{column}{0.48\textwidth}
      \begin{figure}
        \centering
        \fbox{\parbox{0.85\textwidth}{\centering\vspace{2cm}Job 生命周期序列图\\（请替换为 job\_sequence.pdf）\vspace{2cm}}}
        \caption{异步计算任务提交}
      \end{figure}
    \end{column}
  \end{columns}
\end{frame}

%==============================================================================
\section{协议栈详解}
%==============================================================================

\begin{frame}
  \frametitle{链路帧 v2 线缆格式}
  \begin{columns}[T]
    \begin{column}{0.58\textwidth}
      \begin{table}
        \centering
        \scriptsize
        \begin{tabular}{|c|c|p{2.2cm}|p{4cm}|}
          \hline
          \textbf{偏移} & \textbf{长度} & \textbf{字段} & \textbf{说明} \\
          \hline
          0 & 2B & magic & 固定 \texttt{'M'}\texttt{'H'} \\
          \hline
          2 & 1B & version & 当前为 2 \\
          \hline
          3 & 1B & hdr\_len & 固定 19，支持未来扩展 \\
          \hline
          4 & 1B & type & LINK/CTRL/DATA 类型 \\
          \hline
          5 & 1B & flags & 分片/压缩/加密预留 \\
          \hline
          6 & 1B & src & 源 mesh 短地址 \\
          \hline
          7 & 1B & dst & 目的短地址，0xFF 表示未分配/广播 \\
          \hline
          8 & 1B & ttl & 多跳防环 \\
          \hline
          9 & 2B & seq & 发送序号，小端序 \\
          \hline
          11 & 2B & ack & 确认序号，第一版未启用 \\
          \hline
          13 & 2B & payload\_len & 最大 512 \\
          \hline
          15 & 2B & hdr\_crc & 覆盖头部 magic\textasciitilde payload\_len \\
          \hline
          17 & 2B & payload\_crc & 空 payload 时为 0xFFFF \\
          \hline
          19 & N & payload & 上层数据 \\
          \hline
        \end{tabular}
        \caption{link frame v2 格式（总长度 19+N，最大 531B）}
      \end{table}
    \end{column}
    \begin{column}{0.38\textwidth}
      \begin{block}{设计要点}
        \scriptsize
        \begin{itemize}
          \item 双层 CRC：先校验头部，再校验 payload
          \item 小端序，与 HAL 无关
          \item 特殊地址：0x00=Host，0xFF=未分配
          \item 数据面类型：\texttt{0x80} mini9P，\texttt{0x81} RPC，\texttt{0x82} Job，\texttt{0x83} Bulk
        \end{itemize}
      \end{block}
    \end{column}
  \end{columns}
\end{frame}

\begin{frame}
  \frametitle{链路帧类型与解析器}
  \begin{columns}[T]
    \begin{column}{0.48\textwidth}
      \begin{block}{帧类型分类}
        \scriptsize
        \begin{itemize}
          \item \textbf{链路维护} \texttt{0x01\textasciitilde 0x04}：HELLO、HELLO\_ACK、HEARTBEAT、ERROR
          \item \textbf{控制面} \texttt{0x10\textasciitilde 0x1F}：REGISTER、ADDR\_ASSIGN、LEASE\_RENEW/ACK、LINK\_STATE、ROUTE\_UPDATE、HOST\_ADVERTISE、TIME\_SYNC
          \item \textbf{数据面} \texttt{0x80\textasciitilde 0x83}：mini9P、RPC、Job、Bulk
        \end{itemize}
      \end{block}
    \end{column}
    \begin{column}{0.48\textwidth}
      \begin{block}{流式解析器}
        \scriptsize
        \begin{itemize}
          \item 从任意字节流中恢复完整帧
          \item 先寻找 magic 前缀
          \item 收齐 19B 头部并校验 hdr\_crc
          \item 按 payload\_len 收齐 payload 后校验 payload\_crc
          \item 错误时保留尾部 magic 快速重同步
        \end{itemize}
      \end{block}
    \end{column}
  \end{columns}
  \vspace{0.2cm}
  \begin{figure}
    \centering
    \fbox{\parbox{0.85\textwidth}{\centering\vspace{0.8cm}链路解析器状态机图（请替换为 link\_parser\_fsm.pdf）\vspace{0.8cm}}}
    \caption{流式帧解析状态机}
  \end{figure}
\end{frame}

\begin{frame}
  \frametitle{控制面协议：节点注册与地址分配}
  \begin{columns}[T]
    \begin{column}{0.48\textwidth}
      \begin{table}
        \centering
        \scriptsize
        \begin{tabular}{|c|c|p{3.5cm}|}
          \hline
          \textbf{偏移} & \textbf{长度} & \textbf{NODE\_REGISTER} \\
          \hline
          0 & 12B & uid[3]（96-bit 唯一标识） \\
          \hline
          12 & 4B & boot\_id \\
          \hline
          16 & 4B & caps（RELAY=0x1, COORD=0x2） \\
          \hline
          20 & 1B & upstream\_port \\
          \hline
          21 & 3B & 保留，共 24B \\
          \hline
        \end{tabular}
        \caption{节点注册请求（24B）}
      \end{table}
    \end{column}
    \begin{column}{0.48\textwidth}
      \begin{table}
        \centering
        \scriptsize
        \begin{tabular}{|c|c|p{3.5cm}|}
          \hline
          \textbf{偏移} & \textbf{长度} & \textbf{ADDR\_ASSIGN} \\
          \hline
          0 & 12B & uid[3] \\
          \hline
          12 & 4B & boot\_id \\
          \hline
          16 & 4B & lease\_epoch \\
          \hline
          20 & 4B & lease\_ms \\
          \hline
          24 & 1B & addr \\
          \hline
          25 & 1B & flags（OK=0x01） \\
          \hline
          26 & 2B & 保留，共 28B \\
          \hline
        \end{tabular}
        \caption{地址分配响应（28B）}
      \end{table}
    \end{column}
  \end{columns}
  \vspace{0.3cm}
  \begin{alertblock}{关键设计}
    \centering
    \small
    UID + boot\_id 双标识：UID 区分硬件，boot\_id 区分一次启动，避免节点重启后旧租约/路由混淆。
  \end{alertblock}
\end{frame}

\begin{frame}
  \frametitle{控制面协议：租约、链路与路由}
  \begin{columns}[T]
    \begin{column}{0.32\textwidth}
      \begin{table}
        \centering
        \tiny
        \begin{tabular}{|c|c|p{1.6cm}|}
          \hline
          \textbf{偏移} & \textbf{长度} & \textbf{LEASE\_RENEW} \\
          \hline
          0 & 12B & uid \\
          \hline
          12 & 4B & boot\_id \\
          \hline
          16 & 4B & lease\_epoch \\
          \hline
          20 & 1B & addr \\
          \hline
          21 & 3B & 保留 \\
          \hline
        \end{tabular}
        \caption{租约续期 24B}
      \end{table}
    \end{column}
    \begin{column}{0.32\textwidth}
      \begin{table}
        \centering
        \tiny
        \begin{tabular}{|c|c|p{1.6cm}|}
          \hline
          \textbf{偏移} & \textbf{长度} & \textbf{LINK\_STATE} \\
          \hline
          0 & 12B & peer\_uid \\
          \hline
          12 & 4B & peer\_boot\_id \\
          \hline
          16 & 2B & metric \\
          \hline
          18 & 1B & local\_addr \\
          \hline
          19 & 1B & local\_port \\
          \hline
          20 & 1B & peer\_addr \\
          \hline
          21 & 1B & peer\_port \\
          \hline
          22 & 1B & flags \\
          \hline
          23 & 1B & 保留 \\
          \hline
        \end{tabular}
        \caption{链路状态 24B}
      \end{table}
    \end{column}
    \begin{column}{0.32\textwidth}
      \begin{table}
        \centering
        \tiny
        \begin{tabular}{|c|c|p{1.6cm}|}
          \hline
          \textbf{偏移} & \textbf{长度} & \textbf{ROUTE\_UPDATE} \\
          \hline
          0 & 4B & route\_version \\
          \hline
          4 & 2B & metric \\
          \hline
          6 & 1B & dst \\
          \hline
          7 & 1B & next\_hop \\
          \hline
          8 & 1B & action（SET=1, DEL=2） \\
          \hline
          9 & 3B & 保留，共 12B \\
          \hline
        \end{tabular}
        \caption{路由更新 12B}
      \end{table}
    \end{column}
  \end{columns}
  \vspace{0.3cm}
  \begin{block}{路由计算示例}
    \small
    节点 A 上报 ``端口 1 连接到 B'' 后，coordinator 同时生成两条 \texttt{ROUTE\_UPDATE}：
    \begin{itemize}
      \item 告诉 A：dst=B, next\_hop=B，metric=1
      \item 告诉 B：dst=A, next\_hop=A，metric=1
    \end{itemize}
    节点只按 \texttt{route\_version} 更新本地表，旧版本直接忽略。
  \end{block}
\end{frame}

\begin{frame}
  \frametitle{mini9P 协议帧}
  \begin{columns}[T]
    \begin{column}{0.48\textwidth}
      \begin{table}
        \centering
        \scriptsize
        \begin{tabular}{|c|c|p{3.5cm}|}
          \hline
          \textbf{偏移} & \textbf{长度} & \textbf{字段} \\
          \hline
          0 & 1B & version（=1） \\
          \hline
          1 & 1B & type \\
          \hline
          2 & 2B & tag \\
          \hline
          4 & N & payload \\
          \hline
        \end{tabular}
        \caption{mini9P 帧头 4B + payload}
      \end{table}
      \begin{block}{消息类型}
        \scriptsize
        \begin{itemize}
          \item TATTACH/RATTACH、TWALK/RWALK、TOPEN/ROPEN
          \item TREAD/RREAD、TWRITE/RWRITE、TSTAT/RSTAT
          \item TCLUNK/RCLUNK、RERROR（\texttt{0xFF}）
        \end{itemize}
      \end{block}
    \end{column}
    \begin{column}{0.48\textwidth}
      \begin{block}{qid 与 stat 标志}
        \scriptsize
        \begin{itemize}
          \item \texttt{M9P\_QID\_DIR=0x80}：目录
          \item \texttt{M9P\_QID\_VIRTUAL=0x08}：虚拟文件
          \item \texttt{M9P\_QID\_DEVICE=0x04}：设备文件
          \item \texttt{M9P\_QID\_COMPUTE=0x02}：计算文件
          \item \texttt{M9P\_QID\_READONLY=0x01}：只读
        \end{itemize}
      \end{block}
      \begin{block}{限制}
        \scriptsize
        最大路径 255B，文件名 64B，错误文本 64B，帧开销 10B。
      \end{block}
    \end{column}
  \end{columns}
\end{frame}

\begin{frame}
  \frametitle{mini9P 请求/响应示例}
  \begin{table}
    \centering
    \scriptsize
    \begin{tabular}{|c|p{4cm}|p{4cm}|p{4cm}|}
      \hline
      \textbf{类型} & \textbf{关键字段} & \textbf{发送方} & \textbf{用途} \\
      \hline
      TATTACH & fid, requested\_msize, requested\_inflight & 主机 & 建立 session \\
      \hline
      RATTACH & negotiated\_msize, max\_fids, max\_inflight, root\_qid & 节点 & 协商参数 \\
      \hline
      TWALK & fid, newfid, path & 主机 & 定位文件 \\
      \hline
      RWALK & qid & 节点 & 返回对象标识 \\
      \hline
      TREAD & fid, offset, count & 主机 & 读取数据 \\
      \hline
      RREAD & data & 节点 & 返回读取内容 \\
      \hline
      TWRITE & fid, offset, count, data & 主机 & 写入数据 \\
      \hline
      RWRITE & count & 节点 & 返回实际写入长度 \\
      \hline
      RERROR & code, msg & 节点 & 错误响应 \\
      \hline
    \end{tabular}
    \caption{mini9P 主要请求/响应}
  \end{table}
\end{frame}

\begin{frame}
  \frametitle{MCU RPC 协议}
  \begin{columns}[T]
    \begin{column}{0.52\textwidth}
      \begin{table}
        \centering
        \scriptsize
        \begin{tabular}{|c|c|p{3.5cm}|}
          \hline
          \textbf{偏移} & \textbf{长度} & \textbf{字段} \\
          \hline
          0 & 1B & version（=1） \\
          \hline
          1 & 1B & kind（REQUEST/RESPONSE/CANCEL/STREAM\_CHUNK/STREAM\_END） \\
          \hline
          2 & 1B & flags（ONEWAY=0x01, STREAM=0x02） \\
          \hline
          3 & 1B & service\_len \\
          \hline
          4 & 1B & method\_len \\
          \hline
          5 & 1B & 保留 \\
          \hline
          6 & 2B & call\_id \\
          \hline
          8 & 2B & status（OK/NOT\_FOUND/BUSY/...） \\
          \hline
          10 & 4B & deadline\_ms \\
          \hline
          14 & 2B & payload\_len \\
          \hline
          16 & N & service + method + payload \\
          \hline
        \end{tabular}
        \caption{RPC 帧头 16B（最大 512B）}
      \end{table}
    \end{column}
    \begin{column}{0.44\textwidth}
      \begin{block}{特点}
        \scriptsize
        \begin{itemize}
          \item service/method 最大各 31B
          \item STREAM\_CHUNK 的 status 字段复用为 chunk 序号
          \item 支持 oneway 调用
          \item 支持 deadline 超时
          \item 完整性由外层 link payload CRC 保证
        \end{itemize}
      \end{block}
    \end{column}
  \end{columns}
\end{frame}

\begin{frame}
  \frametitle{Job 协议帧}
  \begin{columns}[T]
    \begin{column}{0.55\textwidth}
      \begin{table}
        \centering
        \scriptsize
        \begin{tabular}{|c|c|p{4cm}|}
          \hline
          \textbf{偏移} & \textbf{长度} & \textbf{字段} \\
          \hline
          0 & 1B & version（=1） \\
          \hline
          1 & 1B & kind \\
          \hline
          2 & 1B & state \\
          \hline
          3 & 1B & kernel \\
          \hline
          4 & 2B & request\_id \\
          \hline
          6 & 2B & status \\
          \hline
          8 & 4B & job\_id \\
          \hline
          12 & 2B & progress\_permille \\
          \hline
          14 & 4B & result\_len \\
          \hline
          18 & 2B & payload\_len \\
          \hline
          20 & N & payload（最大 492B） \\
          \hline
        \end{tabular}
        \caption{Job 帧头 20B（最大 512B）}
      \end{table}
    \end{column}
    \begin{column}{0.41\textwidth}
      \begin{block}{消息种类}
        \scriptsize
        \begin{itemize}
          \item CAPS\_REQUEST/RESPONSE
          \item SUBMIT/SUBMIT\_ACK
          \item STATUS\_REQUEST/RESPONSE
          \item RESULT\_REQUEST/RESPONSE
          \item CANCEL\_REQUEST/CANCEL\_ACK
        \end{itemize}
      \end{block}
      \begin{block}{状态}
        \scriptsize
        EMPTY $\rightarrow$ CREATED $\rightarrow$ QUEUED $\rightarrow$ ASSIGNED $\rightarrow$ RUNNING $\rightarrow$ DONE/FAILED/CANCELLED/LOST
      \end{block}
    \end{column}
  \end{columns}
\end{frame}

\begin{frame}
  \frametitle{Host RPC 与 Leader 选举}
  \begin{columns}[T]
    \begin{column}{0.48\textwidth}
      \begin{block}{Host RPC 帧}
        \scriptsize
        \begin{itemize}
          \item TCP/9909，4B 网络序 CBOR body 长度前缀
          \item 最大 payload 1024B，最大帧 1280B
          \item kind：REQUEST/RESPONSE/CANCEL/STREAM\_CHUNK/STREAM\_END
          \item status 包含 \texttt{NOT\_LEADER}，便于客户端重定向
        \end{itemize}
      \end{block}
    \end{column}
    \begin{column}{0.48\textwidth}
      \begin{block}{Leader 选举}
        \scriptsize
        \begin{itemize}
          \item mDNS 发布 \texttt{\_pwos.\_tcp}
          \item 每个主机维护最多 8 个 peer
          \item 比较键：(epoch, priority, host\_uid)
          \item 角色：OBSERVER / FOLLOWER / LEADER
          \item 默认超时 15 秒
        \end{itemize}
      \end{block}
    \end{column}
  \end{columns}
  \vspace{0.3cm}
  \begin{figure}
    \centering
    \fbox{\parbox{0.85\textwidth}{\centering\vspace{0.8cm}多主机 leader 选举与全局命名同步图（请替换为 host\_election.pdf）\vspace{0.8cm}}}
    \caption{主机间协作}
  \end{figure}
\end{frame}

%==============================================================================
\section{计算任务系统}
%==============================================================================

\begin{frame}
  \frametitle{为什么需要独立的计算任务系统？}
  \begin{columns}[T]
    \begin{column}{0.48\textwidth}
      \begin{block}{问题}
        \small
        \begin{itemize}
          \item STM32F4 主频 168MHz，RAM 192KB
          \item 矩阵乘、Mandelbrot、光线追踪一个 tile 可能耗时数百毫秒
          \item 如果放在 service\_task 中执行，会阻塞 mini9P/RPC/控制面
        \end{itemize}
      \end{block}
    \end{column}
    \begin{column}{0.48\textwidth}
      \begin{block}{设计目标}
        \small
        \begin{itemize}
          \item 长计算不阻塞通信
          \item 支持提交、查询、取消、获取结果
          \item 资源受限：固定 job 槽位，无动态内存
          \item 可扩展：新增 kernel 只需注册 stepping 函数
        \end{itemize}
      \end{block}
    \end{column}
  \end{columns}
  \vspace{0.3cm}
  \begin{alertblock}{核心思路}
    \centering
    独立的低优先级 \texttt{compute\_task} + 增量 \texttt{pwos\_compute\_worker\_step()}，每次只推进一个最小工作单元，完成后主动让出 CPU。
  \end{alertblock}
\end{frame}

\begin{frame}
  \frametitle{Job 状态机}
  \begin{figure}
    \centering
    \fbox{\parbox{0.9\textwidth}{\centering\vspace{2cm}Job 状态转换图\\EMPTY $\rightarrow$ CREATED $\rightarrow$ QUEUED $\rightarrow$ ASSIGNED $\rightarrow$ RUNNING $\rightarrow$ DONE/FAILED/CANCELLED/LOST\\（请替换为 job\_state\_machine.pdf）\vspace{2cm}}}
    \caption{Job 生命周期状态机}
  \end{figure}
  \begin{itemize}
    \item \textbf{EMPTY}：槽位空闲
    \item \textbf{QUEUED}：已提交，等待 CPU
    \item \textbf{RUNNING}：compute\_worker 正在增量执行
    \item \textbf{DONE/FAILED/CANCELLED}：终止状态，槽位可被复用
  \end{itemize}
\end{frame}

\begin{frame}
  \frametitle{Kernel 目录与输入格式}
  \begin{table}
    \centering
    \scriptsize
    \begin{tabular}{|c|p{3.5cm}|p{3.5cm}|p{4cm}|}
      \hline
      \textbf{Kernel} & \textbf{输入格式} & \textbf{输出格式} & \textbf{说明} \\
      \hline
      HASH & 任意字节串 & 4B FNV-1a hash & 每字节一步 \\
      \hline
      VECTOR\_ADD & le16 count + A[] + B[] & le16 count + C[] & 最大 32 个 int16 \\
      \hline
      MATMUL & rows, inner, cols, A, B & rows, cols, C & 最大 8x8 \\
      \hline
      MANDELBROT & width, height, max\_iter, cx0, cy0, dx, dy & width, height, max\_iter, pixels & 最大 16x16 \\
      \hline
      RAYTRACE\_TILE & pwos\_render\_request\_t 22B & 16B 头 + RGB565 pixels & 最大 16x7 tile \\
      \hline
    \end{tabular}
    \caption{compute\_worker 内置 kernel}
  \end{table}
  \begin{itemize}
    \item \texttt{PWOS\_COMPUTE\_MAX\_JOBS = 4}，\texttt{INPUT\_CAP = 320B}，\texttt{RESULT\_CAP = 320B}
    \item 每个 job 包含 \texttt{work\_index / work\_total}，用于计算 \texttt{progress\_permille}
  \end{itemize}
\end{frame}

\begin{frame}
  \frametitle{compute\_worker 数据结构}
  \begin{columns}[T]
    \begin{column}{0.58\textwidth}
      \begin{lstlisting}[caption={\scriptsize compute\_worker.h 核心结构}]
typedef struct {
    uint8_t used;
    uint8_t owner_addr;
    uint8_t kernel;
    uint8_t state;
    uint32_t job_id;
    uint16_t progress_permille;
    uint16_t input_len;
    uint16_t result_len;
    uint32_t sequence;
    uint32_t work_index;
    uint32_t work_total;
    uint32_t accumulator;
    uint8_t input[PWOS_COMPUTE_INPUT_CAP];
    uint8_t result[PWOS_COMPUTE_RESULT_CAP];
    char log[PWOS_COMPUTE_LOG_CAP];
} pwos_compute_job_t;
      \end{lstlisting}
    \end{column}
    \begin{column}{0.38\textwidth}
      \begin{block}{字段含义}
        \scriptsize
        \begin{itemize}
          \item \texttt{used}：槽位是否占用
          \item \texttt{owner\_addr + job\_id}：全局唯一键
          \item \texttt{sequence}：FIFO 顺序
          \item \texttt{work\_index / total}：进度
          \item \texttt{accumulator}：hash 累加器
          \item \texttt{log[48]}：人类可读状态
        \end{itemize}
      \end{block}
    \end{column}
  \end{columns}
\end{frame}

\begin{frame}
  \frametitle{增量调度核心代码}
  \begin{lstlisting}[caption={\scriptsize pwos\_compute\_worker\_step() 主循环}]
int pwos_compute_worker_step(pwos_compute_worker_t *worker) {
    worker_lock(worker);
    // 1) 找正在运行的 job
    for (i = 0; i < PWOS_COMPUTE_MAX_JOBS; ++i)
        if (worker->jobs[i].used &&
            worker->jobs[i].state == PWOS_JOB_STATE_RUNNING) {
            job = &worker->jobs[i]; break;
        }
    // 2) 没有则把最老的 QUEUED 切为 RUNNING
    if (job == NULL) {
        for (i = 0; i < PWOS_COMPUTE_MAX_JOBS; ++i)
            if (worker->jobs[i].used &&
                worker->jobs[i].state == PWOS_JOB_STATE_QUEUED &&
                (job == NULL || worker->jobs[i].sequence < job->sequence))
                job = &worker->jobs[i];
        if (job != NULL) {
            job->state = PWOS_JOB_STATE_RUNNING;
            --worker->stats.queued_jobs;
            ++worker->stats.active_jobs;
            ++worker->stats.started;
        }
    }
    // 3) 执行一个 work unit
    switch (job->kernel) {
        case PWOS_JOB_KERNEL_HASH: step_hash(job); break;
        case PWOS_JOB_KERNEL_MATMUL: step_matmul(job); break;
        ...
    }
    ++worker->stats.steps;
    job->progress_permille =
        (job->work_index * PWOS_JOB_PROGRESS_MAX) / job->work_total;
    if (job->work_index >= job->work_total) {
        job->state = PWOS_JOB_STATE_DONE;
        ++worker->stats.completed;
        --worker->stats.active_jobs;
    }
    worker_unlock(worker);
    return 1;
}
  \end{lstlisting}
\end{frame}

\begin{frame}
  \frametitle{Stepping 函数示例}
  \begin{columns}[T]
    \begin{column}{0.48\textwidth}
      \begin{lstlisting}[caption={\scriptsize hash kernel}]
static void step_hash(pwos_compute_job_t *job) {
    job->accumulator ^=
        job->input[job->work_index];
    job->accumulator *= 16777619u;
    ++job->work_index;
    if (job->work_index == job->work_total)
        pwos_job_put_le32(job->result,
                          job->accumulator);
}
      \end{lstlisting}
    \end{column}
    \begin{column}{0.48\textwidth}
      \begin{lstlisting}[caption={\scriptsize matmul kernel}]
static void step_matmul(pwos_compute_job_t *job) {
    uint8_t inner = job->input[1];
    uint8_t cols  = job->input[2];
    uint32_t row = job->work_index / cols;
    uint32_t col = job->work_index % cols;
    uint32_t a_off = 4u;
    uint32_t b_off = a_off + 2u * job->input[0] * inner;
    int32_t sum = 0;
    for (uint32_t k = 0; k < inner; ++k) {
        int32_t a = pwos_job_get_i16(
            job->input + a_off + 2*(row*inner+k));
        int32_t b = pwos_job_get_i16(
            job->input + b_off + 2*(k*cols+col));
        sum += a * b;
    }
    pwos_job_put_i32(
        job->result + 4u + job->work_index*4u, sum);
    ++job->work_index;
}
      \end{lstlisting}
    \end{column}
  \end{columns}
\end{frame}

\begin{frame}
  \frametitle{Raytrace Tile 分布式渲染示例}
  \begin{columns}[T]
    \begin{column}{0.48\textwidth}
      \begin{block}{渲染请求 22B}
        \scriptsize
        \begin{itemize}
          \item scene\_id, tile\_x, tile\_y, tile\_w, tile\_h
          \item image\_w, image\_h, samples, max\_depth
          \item frame\_id, seed, camera\_phase
          \item 最大 tile：16 $\times$ 7，结果最多 240B
        \end{itemize}
      \end{block}
      \begin{block}{执行流程}
        \scriptsize
        \begin{enumerate}
          \item P4 Lua 调度器把 240x320 画面切分为多个 tile
          \item 每个 tile 作为一个 Job 提交到不同 MCU
          \item MCU 用 smallpt 逐像素渲染
          \item F429 LCD 节点接收 tile 结果并显示
        \end{enumerate}
      \end{block}
    \end{column}
    \begin{column}{0.48\textwidth}
      \begin{lstlisting}[caption={\scriptsize raytrace stepping}]
static void step_raytrace_tile(pwos_compute_job_t *job) {
    pwos_render_request_t request;
    if (pwos_render_decode_request(
            job->input, job->input_len, &request) != 0) {
        job->state = PWOS_JOB_STATE_FAILED;
        return;
    }
    uint16_t pixel = pwos_smallpt_render_pixel(
        &request, job->work_index);
    uint32_t offset = PWOS_RENDER_RESULT_HEADER_LEN
                    + job->work_index * 2u;
    job->result[offset]   = (uint8_t)(pixel & 0xFFu);
    job->result[offset+1] = (uint8_t)(pixel >> 8);
    ++job->work_index;
}
      \end{lstlisting}
    \end{column}
  \end{columns}
\end{frame}

\begin{frame}
  \frametitle{计算任务演示流程}
  \begin{columns}[T]
    \begin{column}{0.48\textwidth}
      \begin{block}{Shell 命令示例}
        \scriptsize
        \begin{enumerate}
          \item \texttt{job caps mcu1}：查询能力
          \item \texttt{job submit mcu1 matmul}：提交矩阵乘
          \item \texttt{job status <id>}：查看进度
          \item \texttt{job result <id>}：获取结果
          \item \texttt{job cancel <id>}：取消任务
        \end{enumerate}
      \end{block}
    \end{column}
    \begin{column}{0.48\textwidth}
      \begin{figure}
        \centering
        \fbox{\parbox{0.85\textwidth}{\centering\vspace{2cm}分布式渲染演示截图\\Job 提交 + tile 结果 + LCD 显示\\（请替换为 raytrace\_demo.png）\vspace{2cm}}}
        \caption{计算任务演示}
      \end{figure}
    \end{column}
  \end{columns}
\end{frame}

\begin{frame}
  \frametitle{任务槽位复用与取消}
  \begin{columns}[T]
    \begin{column}{0.48\textwidth}
      \begin{block}{槽位复用策略}
        \small
        \begin{itemize}
          \item 优先找未使用的槽位
          \item 满时选择最早完成的终止态 job（sequence 最小）
          \item 通过 \texttt{stats.slots\_reused} 统计复用次数
        \end{itemize}
      \end{block}
    \end{column}
    \begin{column}{0.48\textwidth}
      \begin{block}{取消机制}
        \small
        \begin{itemize}
          \item 只能在 QUEUED 或 RUNNING 状态取消
          \item 立即把状态置为 CANCELLED
          \item 更新 queued/active 计数
          \item 后续 result 请求返回 CANCELLED 状态
        \end{itemize}
      \end{block}
    \end{column}
  \end{columns}
  \vspace{0.3cm}
  \begin{lstlisting}[caption={\scriptsize 槽位分配与复用}]
static pwos_compute_job_t *alloc_job_locked(pwos_compute_worker_t *worker) {
    pwos_compute_job_t *oldest = NULL;
    for (i = 0; i < PWOS_COMPUTE_MAX_JOBS; ++i) {
        if (worker->jobs[i].used == 0) return &worker->jobs[i];
        if (terminal_state(worker->jobs[i].state) &&
            (oldest == NULL || worker->jobs[i].sequence < oldest->sequence))
            oldest = &worker->jobs[i];
    }
    if (oldest != NULL) ++worker->stats.slots_reused;
    return oldest;
}
  \end{lstlisting}
\end{frame}

%==============================================================================
\section{关键算法}
%==============================================================================

\begin{frame}
  \frametitle{流式帧解析状态机}
  \begin{columns}[T]
    \begin{column}{0.48\textwidth}
      \begin{block}{核心步骤}
        \small
        \begin{enumerate}
          \item 寻找 magic 前缀
          \item 收齐头部并校验 CRC
          \item 读取 payload 长度
          \item 收齐 payload 后最终解码
          \item 错误时保留尾部 magic 快速恢复
        \end{enumerate}
      \end{block}
      \begin{block}{特点}
        \small
        \begin{itemize}
          \item 时间复杂度 O(n)
          \item 固定缓冲区，无动态内存
          \item 对噪声鲁棒
        \end{itemize}
      \end{block}
    \end{column}
    \begin{column}{0.48\textwidth}
      \begin{figure}
        \centering
        \fbox{\parbox{0.85\textwidth}{\centering\vspace{2cm}解析器状态机图\\（请替换为 parser\_fsm.pdf）\vspace{2cm}}}
        \caption{流式帧解析状态机}
      \end{figure}
    \end{column}
  \end{columns}
\end{frame}

\begin{frame}
  \frametitle{地址分配与租约}
  \begin{columns}[T]
    \begin{column}{0.48\textwidth}
      \begin{block}{地址分配}
        \small
        \begin{itemize}
          \item 从 next\_addr 开始线性扫描
          \item 跳过 0x00 和 0xFF
          \item 最多 254 次尝试
          \item 节点地址空间 1\textasciitilde 253
        \end{itemize}
      \end{block}
      \begin{block}{租约机制}
        \small
        \begin{itemize}
          \item 默认租约 30 秒
          \item 节点在租约过半时续期
          \item coordinator 验证 UID + boot\_id
        \end{itemize}
      \end{block}
    \end{column}
    \begin{column}{0.48\textwidth}
      \begin{figure}
        \centering
        \fbox{\parbox{0.85\textwidth}{\centering\vspace{2cm}地址分配与租约时序图\\（请替换为 lease.pdf）\vspace{2cm}}}
        \caption{地址分配与租约}
      \end{figure}
    \end{column}
  \end{columns}
\end{frame}

\begin{frame}
  \frametitle{路由计算}
  \begin{columns}[T]
    \begin{column}{0.48\textwidth}
      \begin{block}{输入}
        \small
        节点上报的链路状态：本地地址、端口、对端地址、metric 等。
      \end{block}
      \begin{block}{输出}
        \small
        收到链路状态后同时生成：
        \begin{itemize}
          \item 正向路由：告诉 local 怎么到 peer
          \item 反向路由：告诉 peer 怎么到 local
        \end{itemize}
      \end{block}
      \begin{block}{节点更新}
        \small
        节点按 route\_version 更新本地表，旧版本忽略。
      \end{block}
    \end{column}
    \begin{column}{0.48\textwidth}
      \begin{figure}
        \centering
        \fbox{\parbox{0.85\textwidth}{\centering\vspace{2cm}路由计算示例图\\（请替换为 route\_calc.pdf）\vspace{2cm}}}
        \caption{路由计算}
      \end{figure}
    \end{column}
  \end{columns}
\end{frame}

\begin{frame}
  \frametitle{typed pending 匹配}
  \begin{columns}[T]
    \begin{column}{0.48\textwidth}
      \begin{block}{问题}
        \small
        主机同时向一个节点发送文件、调用、任务请求，如何正确匹配响应？
      \end{block}
      \begin{block}{匹配键}
        \small
        \begin{itemize}
          \item 数据类型（mini9P/RPC/Job）
          \item 源地址
          \item 标签（wire\_tag / call\_id / request\_id）
        \end{itemize}
      \end{block}
      \begin{block}{流式响应}
        \small
        按 chunk 序号顺序追加，不连续则报错，STREAM\_END 才唤醒。
      \end{block}
    \end{column}
    \begin{column}{0.48\textwidth}
      \begin{figure}
        \centering
        \fbox{\parbox{0.85\textwidth}{\centering\vspace{2cm}pending 匹配流程图\\（请替换为 pending\_match.pdf）\vspace{2cm}}}
        \caption{typed pending 匹配}
      \end{figure}
    \end{column}
  \end{columns}
\end{frame}

\begin{frame}
  \frametitle{增量计算调度}
  \begin{columns}[T]
    \begin{column}{0.48\textwidth}
      \begin{block}{问题}
        \small
        MCU 计算能力弱，长计算不能阻塞通信和控制面。
      \end{block}
      \begin{block}{方案}
        \small
        \begin{itemize}
          \item 低优先级 compute\_task
          \item 每次只执行一个 work unit
          \item 完成后主动让出 CPU
          \item 支持进度查询
        \end{itemize}
      \end{block}
    \end{column}
    \begin{column}{0.48\textwidth}
      \begin{figure}
        \centering
        \fbox{\parbox{0.85\textwidth}{\centering\vspace{2cm}增量调度示意图\\（请替换为 incremental\_compute.pdf）\vspace{2cm}}}
        \caption{增量计算调度}
      \end{figure}
    \end{column}
  \end{columns}
\end{frame}

%==============================================================================
\section{代码实现与数据流}
%==============================================================================

\begin{frame}
  \frametitle{关键代码片段}
  \begin{columns}[T]
    \begin{column}{0.48\textwidth}
      \begin{block}{片段 1：coordinator 主循环帧分发}
        \scriptsize
        \texttt{handle\_frame()} 根据帧类型分发：HELLO 回复 ACK，注册/租约/链路状态进入对应 handler，数据面交给 session\_manager 匹配 pending。
      \end{block}
      \begin{block}{片段 2：地址分配线性扫描}
        \scriptsize
        \texttt{alloc\_addr()} 从 \texttt{next\_addr} 开始扫描 1\textasciitilde 253，跳过已用地址，最多 254 次尝试。
      \end{block}
    \end{column}
    \begin{column}{0.48\textwidth}
      \begin{block}{片段 3：pending 匹配}
        \scriptsize
        \texttt{session\_manager\_deliver\_data()} 遍历 pending 表，严格匹配数据类型、源地址、标签，并校验 boot\_id 是否一致。
      \end{block}
      \begin{block}{片段 4：compute\_worker 增量 step}
        \scriptsize
        \texttt{pwos\_compute\_worker\_step()} 每次推进一个 work unit，更新进度，完成后让出 CPU。
      \end{block}
    \end{column}
  \end{columns}
  \vspace{0.3cm}
  \begin{figure}
    \centering
    \fbox{\parbox{0.6\textwidth}{\centering\vspace{0.6cm}关键代码截图（请替换为 code\_snippet.png）\vspace{0.6cm}}}
    \caption{关键代码片段}
  \end{figure}
\end{frame}

\begin{frame}
  \frametitle{完整数据流}
  \begin{figure}
    \centering
    \fbox{\parbox{0.9\textwidth}{\centering\vspace{2.5cm}主机到 STM32 的完整调用链图\\包含 mini9P read、Job 提交两条主线\\（请替换为 dataflow.pdf）\vspace{2.5cm}}}
    \caption{完整数据流与调用链}
  \end{figure}
  \begin{itemize}
    \item 主机侧：路径解析 $\rightarrow$ session $\rightarrow$ UART 发送
    \item 节点侧：UART ISR $\rightarrow$ 任务队列 $\rightarrow$ node\_control $\rightarrow$ service\_runtime / compute\_task $\rightarrow$ 响应返回
  \end{itemize}
\end{frame}

\begin{frame}
  \frametitle{调用链：\texttt{cat /mcu1/sys/health}}
  \begin{table}
    \centering
    \scriptsize
    \begin{tabular}{|c|p{4cm}|p{6cm}|}
      \hline
      \textbf{层级} & \textbf{主机侧} & \textbf{节点侧} \\
      \hline
      用户入口 & WebShell / HTTP & -- \\
      \hline
      命名空间 & cluster\_vfs 解析 /mcu1/sys/health & local\_vfs 生成 /sys/health \\
      \hline
      会话 & session\_manager 获取/创建 client & service\_runtime 分发 \\
      \hline
          协议 & mini9P client：TATTACH/TWALK/TOPEN/TREAD & mini9P server 处理 \\
      \hline
      传输 & link frame DATA\_MINI9P & node\_control 转发/投递 \\
      \hline
      物理 & UART TX & UART RX DMA + IDLE 中断 \\
      \hline
    \end{tabular}
    \caption{读取健康状态的完整调用链}
  \end{table}
\end{frame}

\begin{frame}
  \frametitle{调用链：\texttt{job submit mcu1 matmul}}
  \begin{table}
    \centering
    \scriptsize
    \begin{tabular}{|c|p{4cm}|p{6cm}|}
      \hline
      \textbf{层级} & \textbf{主机侧} & \textbf{节点侧} \\
      \hline
      用户入口 & host\_shell：job submit & -- \\
      \hline
      任务管理 & job\_manager 分配 host\_job\_id & job\_service 接收 \\
      \hline
      会话 & session\_manager 发送 DATA\_JOB SUBMIT & service\_runtime 解析 \\
      \hline
      协议 & Job 协议：SUBMIT / STATUS / RESULT & Job 协议响应 \\
      \hline
      计算 & -- & compute\_worker 入队、compute\_task 增量执行 \\
      \hline
      结果 & job\_result 命令轮询 & 返回 RESULT\_RESPONSE \\
      \hline
    \end{tabular}
    \caption{提交计算任务的完整调用链}
  \end{table}
\end{frame}

%==============================================================================
\section{错误处理与容错}
%==============================================================================

\begin{frame}
  \frametitle{CRC 与噪声恢复}
  \begin{columns}[T]
    \begin{column}{0.48\textwidth}
      \begin{block}{双层 CRC}
        \small
        \begin{itemize}
          \item 头部 CRC 覆盖 magic\textasciitilde payload\_len
          \item payload CRC 覆盖 payload，空 payload 为 0xFFFF
          \item 头部错误可立即丢弃，不必等待完整帧
        \end{itemize}
      \end{block}
    \end{column}
    \begin{column}{0.48\textwidth}
      \begin{block}{重同步策略}
        \small
        \begin{itemize}
          \item 解码失败时扫描缓冲区中的下一个 magic
          \item 保留未处理字节，避免丢整段数据
          \item 流式 parser 支持任意分片
        \end{itemize}
      \end{block}
    \end{column}
  \end{columns}
  \vspace{0.3cm}
  \begin{figure}
    \centering
    \fbox{\parbox{0.85\textwidth}{\centering\vspace{0.8cm}噪声帧恢复示意图（请替换为 noise\_recovery.pdf）\vspace{0.8cm}}}
    \caption{链路层错误恢复}
  \end{figure}
\end{frame}

\begin{frame}
  \frametitle{超时、重试与租约过期}
  \begin{columns}[T]
    \begin{column}{0.48\textwidth}
      \begin{block}{请求超时}
        \small
        \begin{itemize}
          \item RPC 带 deadline\_ms
          \item session\_manager pending 槽带超时
          \item 超时后清理 pending，返回 DEADLINE
        \end{itemize}
      \end{block}
    \end{column}
    \begin{column}{0.48\textwidth}
      \begin{block}{租约过期}
        \small
        \begin{itemize}
          \item 节点在租约过半续期
          \item coordinator 超时未收到续期则释放地址
          \item 节点重启后 boot\_id 变化，旧状态不混淆
        \end{itemize}
      \end{block}
    \end{column}
  \end{columns}
  \vspace{0.3cm}
  \begin{block}{节点丢失检测}
    \small
    结合 HELLO 超时、LEASE 超时、LINK\_STATE 缺失，coordinator 标记节点 LOST，清理路由和 pending。
  \end{block}
\end{frame}

\begin{frame}
  \frametitle{Job 取消与失败处理}
  \begin{columns}[T]
    \begin{column}{0.48\textwidth}
      \begin{block}{取消机制}
        \small
        \begin{itemize}
          \item CANCEL\_REQUEST 在 QUEUED/RUNNING 时生效
          \item 立即转 CANCELLED，更新计数
          \item result 请求返回 CANCELLED 状态
        \end{itemize}
      \end{block}
    \end{column}
    \begin{column}{0.48\textwidth}
      \begin{block}{失败来源}
        \small
        \begin{itemize}
          \item 输入校验失败：BAD\_REQUEST
          \item 槽位满：BUSY
          \item 未知 kernel：UNSUPPORTED
          \item 渲染请求解码失败：FAILED
          \item 结果未准备好：NOT\_READY
        \end{itemize}
      \end{block}
    \end{column}
  \end{columns}
\end{frame}

%==============================================================================
\section{工程亮点与实验验证}
%==============================================================================

\begin{frame}
  \frametitle{工程设计亮点}
  \begin{columns}[T]
    \begin{column}{0.48\textwidth}
      \begin{block}{架构层面}
        \small
        \begin{itemize}
          \item 控制面与数据面严格分离
          \item 主机集中协调，节点只维护本地状态
          \item 多主机 leader 选举与全局命名
          \item 统一文件抽象覆盖诊断、RPC、Job
        \end{itemize}
      \end{block}
    \end{column}
    \begin{column}{0.48\textwidth}
      \begin{block}{实现层面}
        \small
        \begin{itemize}
          \item STM32 端零动态内存
          \item 固定任务、固定帧池、固定 job 槽位
          \item typed pending 精确匹配
          \item generation 机制保证一致性
          \item 共享协议模块提供 PC 单元测试
        \end{itemize}
      \end{block}
    \end{column}
  \end{columns}
\end{frame}

\begin{frame}
  \frametitle{实验验证}
  \begin{columns}[T]
    \begin{column}{0.48\textwidth}
      \begin{block}{单主机闭环}
        \small
        \begin{itemize}
          \item \texttt{cat /mcu1/sys/health}
          \item \texttt{rpc mcu1 system.ping hello}
          \item \texttt{job submit mcu1 matmul}
          \item \texttt{job submit mcu1 mandelbrot}
        \end{itemize}
      \end{block}
    \end{column}
    \begin{column}{0.48\textwidth}
      \begin{block}{多主机闭环}
        \small
        \begin{itemize}
          \item P4 与 S3 接入同一 LAN
          \item 双方看到相同 leader
          \item 跨 owner 读取 \texttt{/mcuN/sys/health}
        \end{itemize}
      \end{block}
    \end{column}
  \end{columns}
  \vspace{0.3cm}
  \begin{figure}
    \centering
    \fbox{\parbox{0.85\textwidth}{\centering\vspace{1cm}实验截图/串口日志/WebShell 截图\\（请替换为 experiment.png）\vspace{1cm}}}
    \caption{实验验证截图}
  \end{figure}
\end{frame}

\begin{frame}
  \frametitle{性能与可扩展性}
  \begin{columns}[T]
    \begin{column}{0.48\textwidth}
      \begin{block}{性能特点}
        \small
        \begin{itemize}
          \item 帧解析：O(n)，固定缓冲区
          \item 地址分配：最多 254 次扫描
          \item 路由查找：小表线性扫描
          \item 计算任务：增量执行，不阻塞通信
          \item 结果进度粒度：1\textperthousand
        \end{itemize}
      \end{block}
    \end{column}
    \begin{column}{0.48\textwidth}
      \begin{block}{可扩展方向}
        \small
        \begin{itemize}
          \item 扩展地址空间
          \item 引入加权路由 metric
          \item 实现大数据分片传输
          \item 增加主机间安全机制
        \end{itemize}
      \end{block}
    \end{column}
  \end{columns}
  \vspace{0.3cm}
  \begin{figure}
    \centering
    \fbox{\parbox{0.6\textwidth}{\centering\vspace{0.6cm}性能测试图表（请替换为 performance.pdf）\vspace{0.6cm}}}
    \caption{性能与扩展性}
  \end{figure}
\end{frame}

%==============================================================================
\section{不足与改进}
%==============================================================================

\begin{frame}
  \frametitle{当前不足与改进方向}
  \begin{table}
    \centering
    \small
    \begin{tabular}{|p{0.4\textwidth}|p{0.5\textwidth}|}
      \hline
      \textbf{当前不足} & \textbf{改进方向} \\
      \hline
      地址空间 8 位，最多 253 节点 & 扩展为 16 位或分层地址 \\
      \hline
      路由 metric 仅基于跳数 & 引入带宽、负载、延迟加权 \\
      \hline
      大数据传输未实现 & 实现分片传输协议 \\
      \hline
      主机间无安全机制 & 增加预共享密钥或证书 \\
      \hline
      Job 调度无优先级 & 增加优先级队列与抢占点 \\
      \hline
      推理服务仍是骨架 & 完善分布式推理协作 \\
      \hline
    \end{tabular}
    \caption{不足与改进方向}
  \end{table}
\end{frame}

%==============================================================================
\section{总结}
%==============================================================================

\begin{frame}
  \frametitle{总结}
  \begin{itemize}
    \item \textbf{What：}9P-Walkers 是一个面向异构 MCU 集群的实验性分布式系统，核心是统一命名空间下的跨节点资源访问。
    \item \textbf{How：}采用 ESP32 主机 + STM32 节点、自定义链路协议、mini9P 文件抽象、分层任务调度和多主机协作。
    \item \textbf{成果：}实现了节点自动发现、地址租约、路由下发、远程文件/调用/任务、多主机全局命名和上板验证。
    \item \textbf{价值：}为资源受限设备上的集群计算、远程诊断和任务调度提供了一套可扩展、可测试的实现基础。
  \end{itemize}
  \vspace{0.3cm}
  \begin{alertblock}{答辩核心结论}
    \centering
    在受限资源与有限时间下，\alert{统一命名 + 自动组网 + 远程调用 + 分布式计算} 的最小闭环已实现、可验证、可扩展。
  \end{alertblock}
\end{frame}

\begin{frame}
  \frametitle{致谢}
  \centerline{\Large 谢谢！}
  \vspace{0.4cm}
  \centerline{\normalsize 欢迎老师和同学提出问题}
\end{frame}

\end{document}
